% !TeX root = ../main.tex
\section{Efficient simulation in the phase polynomial representation}
\label{appendix:efficient-simulation}

This appendix establishes the runtime and memory bounds for simulation in the phase-polynomial representation.
We analyze the cost of state updates and compatibility tests, then show how to compute amplitudes and Pauli expectation values, sample computational-basis outcomes, and test equality of represented states.

\begin{efficientsimulationrestatement}[Efficient simulation of 3PP circuits]
\efficientSimulationStatement{}
\end{efficientsimulationrestatement}

\begin{proof}
We first establish the cost of updating the state representation and checking measurement compatibility, then prove the bounds for queries on represented states.
Throughout, we use dense coefficient arrays and ordinary Gaussian elimination, with arithmetic on phase coefficients performed modulo $8$.

The affine support is specified by $x_0$ and $B$, requiring $O(n^2)$ bits.
The phase polynomial has the weighted form
\begin{align}
    p(y)
    = d+\sum_i a_i y_i
    +2\sum_{i<j}b_{ij}y_i y_j
    +4\sum_{i<j<k}c_{ijk}y_i y_j y_k,
\end{align}
where $d,a_i\in\mathbb{Z}_8$, $b_{ij}\in\mathbb{Z}_4$, and $c_{ijk}\in\mathbb{F}_2$.
Since $r\leq n$, storing these coefficients requires $O(n^3)$ bits.

We will repeatedly use two elementary facts about polynomial updates.
First, each substitution $y_i\mapsto y_i\oplus 1$ or $y_i\mapsto y_i\oplus y_j$ can be performed in $O(r^2)$ time, since only the $O(r^2)$ monomials containing $y_i$ are affected, and each produces a bounded number of terms.
Second, an affine Boolean function can be expanded as a weighted polynomial modulo $8$ in $O(r^3)$ time.
Indeed, for a parity of variables,
\begin{align}
    \bigoplus_{i\in I}y_i
    =\sum_{i\in I}y_i
    -2\sum_{\substack{i<j\\i,j\in I}}y_i y_j
    +4\sum_{\substack{i<j<k\\i,j,k\in I}}y_i y_j y_k
    \pmod 8,
\end{align}
and an affine constant is incorporated using $1\oplus f=1-f$.

We now consider circuit evolution.
Preparations append a physical coordinate and, for a $\ket{+}$ preparation, a free support coordinate.
An $X$ gate changes $x_0$, while a $\CNOT$ gate applies the corresponding row operation to $x_0$ and $B$.
Neither requires changing $p$.
A $\T$ or $Z$ gate adds an integer multiple of the affine Boolean function describing the affected physical bit to $p$.
These updates, and hence arbitrary Pauli updates, take at most $O(n^3)$ time.

For a $Z$-type measurement, write the measured operator as $Z(v)$.
Its value on a support point is determined by the affine function
\begin{align}
    f(y)=\langle v,x_0\oplus By\rangle.
\end{align}
If $f$ is constant, the outcome is deterministic and the state is unchanged.
Otherwise, each outcome has probability $1/2$ and imposes one affine equation on $y$.
Choose a variable with nonzero coefficient in that equation and use $O(r)$ elementary affine substitutions to make the equation a constraint on that variable alone.
Applying these substitutions to $p$ costs $O(r^3)$, and updating the support parametrization costs at most $O(nr^2)$.
Fixing the constrained variable and removing it therefore gives the conditional-state representation in $O(n^3)$ time.

A measurement of a nonidentity $X$-type Pauli string of weight $w$ reduces to a single-qubit measurement by \cref{eq:x-measurement-reduction}, using at most $2(w-1)$ CNOTs in total.
Each CNOT updates rows of $x_0$ and $B$ in $O(n)$ time and leaves $p$ unchanged, so the reduction adds only $O(n^2)$ time.
The $O(n^3)$ expected cost established below therefore applies to one original Pauli-string measurement, without changing the operation count $m$ or the total runtime.
The identity measurement is deterministic and leaves the state unchanged.

Next, consider a proposed measurement of $X_a$.
Gaussian elimination decides whether $e_a\in\operatorname{im}B$ in $O(n^3)$ time.
If not, the measurement is compatible case~(1).
For outcome $\mu$, append $e_a$ as a new column of $B$ and replace $p$ by
\begin{align}
    p'(y,z)=p(y)+4\mu z.
\end{align}
Both outcomes have probability $1/2$, and the update takes at most $O(n^3)$ time.

If $e_a\in\operatorname{im}B$, solve $Bt=e_a$ and construct
\begin{align}
    D(y)=p(y\oplus t)-p(y).
\end{align}
The translation $y\mapsto y\oplus t$ uses at most $r$ elementary substitutions, so constructing $D$ costs $O(r^3)$.
By the discrete-derivative property, $D$ is a Clifford phase polynomial and therefore has only $O(r^2)$ coefficients.

We now check the three internal-measurement cases and bound their updates.
For case~(2a), coefficient inspection determines whether $D=4f$ for an affine Boolean function $f$.
Conditioning on $f(y)=\mu$ is then the same affine restriction used for $Z$ measurements, costing $O(n^3)$ time.

For case~(2b), coefficient inspection first determines whether $D=2+4g$ for a quadratic Boolean function $g$.
The compatibility criterion is that the alternating matrix of its quadratic part has rank at most two, which can be checked in $O(r^3)$ time.
When this condition holds, binary elimination produces affine $\ell$ and linear $u,v$ such that
\begin{align}
    g(y)=\ell(y)\oplus u(y)v(y).
\end{align}
Here one may take $u=v=0$ when the quadratic part vanishes.
Using
\begin{align}
    2g=2\ell+2uv-4\ell uv\pmod 8,
\end{align}
we can construct the weighted coefficients of $2g$ in $O(r^3)$ time.
For the term $2uv$, only the linear and quadratic parts of the parity expansions of $u$ and $v$ contribute modulo $8$.
For the term $4\ell uv$, it suffices to multiply the affine forms over $\mathbb{F}_2$.
Thus the conditional phase
\begin{align}
    p'=p-2(-1)^\mu g
\end{align}
can be constructed within the same bound.

For case~(2c), let $d=D(0)$.
If $d$ is odd, the linear coefficients of $D$ determine the only possible affine Boolean function $h$ with $h(0)=0$ satisfying
\begin{align}
    D(y)=d(-1)^{h(y)}.
\end{align}
Comparing coefficients verifies this identity in polynomial time bounded by $O(r^3)$.
If it holds, expanding $h$ and forming
\begin{align}
    p'=p+(4\mu-d)h
\end{align}
also costs $O(r^3)$.
If none of these cases applies, the measurement is incompatible.
Consequently, compatibility checking and every compatible conditional-state update take $O(n^3)$ time.

It remains to account for sampling measurement outcomes.
The possible probabilities belong to the fixed set
\begin{align}
    \left\{0,\frac12,1,\frac{2-\sqrt2}{4},
    \frac{2+\sqrt2}{4}\right\}.
\end{align}
The rational probabilities are sampled directly using fair random bits.
For either irrational probability, compare a uniformly random binary expansion with successively refined rational bounds on that probability.
After $k$ bits and bounds of precision $O(2^{-k})$, the probability that the comparison remains unresolved is $O(2^{-k})$.
Since the bounds can be computed in time polynomial in $k$, this procedure has constant expected runtime.
Combining outcome sampling with the updates above gives $O(n^3)$ expected time per operation and $O(mn^3)+T_{\mathrm{classical}}$ expected time for the circuit.

We next establish the costs of queries on represented states.
For a computational-basis amplitude, Gaussian elimination decides whether
\begin{align}
    By=x\oplus x_0
\end{align}
has a solution in $O(n^3)$ time.
Since $B$ has full column rank, any solution is unique, and
\begin{align}
    \langle x|\psi\rangle
    =
    \begin{cases}
        2^{-r/2}\omega_8^{p(y)}, & By=x\oplus x_0,\\
        0, & x\notin x_0\oplus\operatorname{im}B.
    \end{cases}
\end{align}
Evaluating $p(y)$ costs $O(r^3)$, proving the amplitude bound.
For computational-basis sampling, choose $y$ uniformly in $\mathbb{F}_2^r$ and return $x_0\oplus By$.
This costs $O(nr)=O(n^2)$ time.

To compute a Pauli expectation, write
\begin{align}
    P=i^\kappa X(a)Z(b).
\end{align}
If $a\notin\operatorname{im}B$, the translated support is disjoint from the original support and the expectation is zero.
Otherwise, solve $Bt=a$ and expand
\begin{align}
    \bra{\psi}P\ket{\psi}
    =i^\kappa 2^{-r}
    \sum_{y\in\mathbb{F}_2^r}
    \omega_8^{p(y)-p(y\oplus t)
    +4\langle b,x_0\oplus By\rangle}.
    \label{eq:pauli-expectation-gauss-sum}
\end{align}
Constructing the exponent costs $O(n^3)$ by the translation procedure above.
It is a Clifford phase polynomial, so the remaining sum is a quadratic Gauss sum over $\mathbb{Z}_4$.
Reduction of the quadratic form by binary elimination evaluates this sum exactly in $O(r^3)$ time, giving the claimed $O(n^3)$ bound.
The result can be retained exactly using a root of unity and a power of $\sqrt{2}$, including the possibility of zero.

Finally, consider equality of two represented states.
Gaussian elimination checks whether their affine supports coincide in $O(n^3)$ time.
If they do not, the states are unequal.
If they do, their support coordinates are related by an invertible affine transformation, which can also be found in $O(n^3)$ time.
Decompose this transformation into $O(r^2)$ elementary affine substitutions and apply them to one phase polynomial.
Since each substitution costs $O(r^2)$, expressing both phases in common coordinates costs $O(n^4)$ time.

In these common coordinates, the states are equal up to global phase exactly when their phase difference is constant modulo $8$.
Uniqueness of the multilinear polynomial representation over $\mathbb{Z}_8$ makes this equivalent to the vanishing of every nonconstant coefficient of the difference.
Checking these coefficients costs $O(r^3)$, completing the $O(n^4)$ equality bound.
\end{proof}
